\documentclass[11pt]{article}
\usepackage[letterpaper,margin=0.88in]{geometry}
\usepackage[T1]{fontenc}
\usepackage{lmodern}
\usepackage{microtype}
\usepackage{amsmath,amssymb,mathtools}
\usepackage{booktabs,array}
\usepackage{xcolor}
\usepackage{tcolorbox}
\usepackage{enumitem}
\usepackage{hyperref}
\hypersetup{
  colorlinks=true,
  linkcolor=blue!50!black,
  citecolor=blue!50!black,
  urlcolor=blue!50!black,
  pdftitle={Prediction from Partial Marginals in Graphical Models},
  pdfauthor={Product Model Project}
}
\setlist[itemize]{leftmargin=1.45em,itemsep=3pt,topsep=4pt}
\setlist[enumerate]{leftmargin=1.55em,itemsep=3pt,topsep=4pt}
\newcommand{\X}{\mathbf X}
\newcommand{\E}{\mathbb E}
\newcommand{\Pp}{\mathbb P}
\newcommand{\KL}{D_{\mathrm{KL}}}
\newcommand{\1}{\mathbf 1}
\newcommand{\ind}{\mathop{\perp}\!\!\!\perp}
\newcommand{\mcC}{\mathcal C}
\newcommand{\mcA}{\mathcal A}
\newcommand{\mcZ}{\mathcal Z}
\newcommand{\hatp}{\widehat p}
\newcommand{\softmax}{\operatorname{softmax}}

\definecolor{navy}{HTML}{16324F}
\definecolor{teal}{HTML}{0C7185}
\definecolor{pale}{HTML}{EAF4F6}
\definecolor{rule}{HTML}{9CB9C1}

\newtcolorbox{keyidea}{
  colback=pale,
  colframe=rule,
  boxrule=0.7pt,
  arc=1.5pt,
  left=8pt,right=8pt,top=7pt,bottom=7pt
}
\newtcolorbox{warningbox}{
  colback=yellow!8,
  colframe=orange!55!black,
  boxrule=0.6pt,
  arc=1.5pt,
  left=8pt,right=8pt,top=7pt,bottom=7pt
}

\begin{document}

\begin{center}
{\LARGE\bfseries\color{navy} Prediction from Partial Marginals\\[4pt]
in High-Alphabet Graphical Models}\\[9pt]
{\large A maximum-entropy, Bayesian, message-passing, and compression perspective}\\[8pt]
{\small August 2026}
\end{center}

\vspace{4pt}
\hrule height 0.8pt
\vspace{12pt}

\begin{abstract}
Suppose that the next token \(Y\) and selected past tokens
\(X_1,\ldots,X_k\) take values in a very large alphabet.  A complete conditional
table \(p(Y\mid X_1,\ldots,X_k)\) cannot be estimated.  What may be available are
selected low-order marginals: target--past pairs \(p(Y,X_i)\), past--past pairs
\(p(X_i,X_j)\), and target--past--past triples \(p(Y,X_i,X_j)\), perhaps only for
some lag pairs.  This note gives a concrete way to turn such partial information
into a next-token predictor.

The note develops two complementary completions.  Maximum entropy, or more
generally minimum relative entropy relative to a baseline model, selects one
joint distribution.  Bayesian averaging instead places a prior over possible
joint distributions or factor parameters and averages their predictions.  The
note derives the factor-graph potentials for the first route, explains how
past--past marginals affect the final conditional prediction, and makes the
Bayesian route explicit.  In particular, it gives an exactly conjugate
Dirichlet construction for a fixed tree, its decomposable-graph analogue, and
the reason general loopy models are not cheaply conjugate.  It also studies
uncertainty about which interactions are trustworthy, including a factorized
soft-reliability model and exact averaging over tree structures.  It
distinguishes exact tree and junction-tree calculations from approximate
message passing and ends with a practical path for a causal compressor.
\end{abstract}

\newpage
{\footnotesize\tableofcontents}
\newpage

\section{The statistical problem}

Let
\[
  Y = \text{the next token},\qquad
  \X=(X_1,\ldots,X_k)=\text{selected tokens from the decoded past}.
\]
The variables may be particular lags, such as
\(X_1=x_{t-1}\), \(X_2=x_{t-5}\), \(X_3=x_{t-100}\), or more general
causal features of the already decoded prefix.  Let \(\mcA\) be the token
alphabet, possibly with \(\lvert\mcA\rvert\) around \(10^5\).

For prediction and compression, the desired object is
\[
  q_t(y\mid x_1,\ldots,x_k).
\]
The arithmetic or range coder spends approximately
\[
  -\log_2 q_t(y_t\mid x_{1,t},\ldots,x_{k,t})
\]
bits on token \(y_t\).  Encoder and decoder must construct exactly the same
distribution from the already decoded past.

The difficulty is that a literal conditional table has
\(\lvert\mcA\rvert^{k+1}\) cells.  It cannot be learned from realistic data.
Instead, suppose measurements give only selected low-order distributions such
as
\[
  \hatp_{Yi}(y,a),\qquad
  \hatp_{ij}(a,b),\qquad
  \hatp_{Yij}(y,a,b).
\]
The hats emphasize that these are empirical estimates, not necessarily exact
population quantities.

\begin{keyidea}
Partial marginals do \emph{not} specify a unique joint distribution.  They
specify a set of compatible joint distributions.  A prediction rule requires
an additional modelling principle for choosing one member of that set.
\end{keyidea}

\section{The constraint hypergraph}

It is useful to organize the available information as a hypergraph.  Its
vertices are
\[
  V=\{Y,X_1,\ldots,X_k\},
\]
and its hyperedges are the variable subsets whose marginal distributions are
measured and trusted.

For example, the collection
\[
  \mcC=
  \bigl\{
    \{Y,X_1\},\{Y,X_2\},\{Y,X_5\},
    \{X_1,X_2\},\{X_2,X_5\},
    \{Y,X_1,X_2\},\{Y,X_5,X_6\}
  \bigr\}
\]
means that some target--past pairs, some past--past pairs, and two selected
target--past--past triples are available.  There is no requirement that every
pair or every triple be present.

For a joint distribution \(q\) on all variables, write \(q_C\) for its marginal
on a subset \(C\in\mcC\).  In the idealized exact-constraint problem, the
feasible set is
\[
  \mathcal F
  =
  \left\{
  q:
  q_C=\hatp_C \text{ for every } C\in\mcC
  \right\}.
\]
Before doing anything else, overlapping empirical marginals should be checked
for consistency.  For instance, the one-variable distribution of \(X_i\)
obtained from \(\hatp_{Yi}\) should agree with the one obtained from
\(\hatp_{ij}\).  Finite samples usually violate such identities slightly; this
is addressed in Section~\ref{sec:noisy}.

\section{Maximum entropy: the least-committal completion}

The maximum-entropy completion is
\begin{equation}
  q^\star
  =
  \arg\max_{q\in\mathcal F}
  H(q),
  \qquad
  H(q)=-\sum_{z\in\mcZ}q(z)\log q(z).
  \label{eq:maxent}
\end{equation}
It preserves exactly the specified marginal information and adds no
unconstrained interaction deliberately.  This is the standard maximum-entropy
principle and is equivalent to an exponential-family graphical model
\cite{wainwright_jordan}.

\subsection{Deriving the potentials}

For clarity, assume first that every feasible distribution has positive
probability.  Introduce one multiplier
\(\theta_C(c)\) for every measured configuration \(c\) of every clique
\(C\in\mcC\).  The Lagrangian is
\begin{align}
  \mathcal L(q,\lambda,\theta)
  ={}&
  -\sum_z q(z)\log q(z)
  +\lambda\left(\sum_z q(z)-1\right) \notag\\
  &+\sum_{C\in\mcC}\sum_c
  \theta_C(c)
  \left[
    \sum_{z:z_C=c}q(z)-\hatp_C(c)
  \right].
  \label{eq:lagrangian}
\end{align}
Differentiating with respect to \(q(z)\) and setting the derivative to zero
gives
\[
  -\log q(z)-1+\lambda+\sum_{C\in\mcC}\theta_C(z_C)=0.
\]
Hence
\begin{equation}
  q_\theta(z)
  =
  \frac{1}{Z(\theta)}
  \exp\left[
    \sum_{C\in\mcC}\theta_C(z_C)
  \right],
  \label{eq:factor_model}
\end{equation}
where
\[
  Z(\theta)
  =
  \sum_z
  \exp\left[
    \sum_{C\in\mcC}\theta_C(z_C)
  \right].
\]
The functions \(\theta_C\) are the \emph{potentials}.  They are not generally
\(\log\hatp_C\).  Because cliques overlap, a potential must compensate for all
the other potentials so that every required marginal is correct.

The fitting equations are the moment-matching equations
\begin{equation}
  q_{\theta,C}(c)
  =
  \hatp_C(c)
  \qquad
  \text{for all }C\in\mcC,\ c.
  \label{eq:moment_match}
\end{equation}
Equivalently, define the concave dual objective
\begin{equation}
  \ell(\theta)
  =
  \sum_{C\in\mcC}\sum_c
  \hatp_C(c)\theta_C(c)
  -
  \log Z(\theta).
  \label{eq:dual}
\end{equation}
Its gradient is
\begin{equation}
  \frac{\partial\ell}{\partial\theta_C(c)}
  =
  \hatp_C(c)-q_{\theta,C}(c).
  \label{eq:gradient}
\end{equation}
Thus every fitting method is trying to make the model marginal equal to the
measured marginal.

\begin{warningbox}
Potentials are not unique: constants can be shifted between overlapping
potentials without changing \(q_\theta\).  The probability distribution
\(q_\theta\), and hence its prediction, is the meaningful object.
\end{warningbox}

\section{The full pair-and-selected-triple model}

Let \(E_X\) denote selected past--past pairs, and \(E_Y\) selected
target--past--past triples.  A useful explicit model is
\begin{align}
  \log q(y,x)
  ={}&
  \theta_Y(y)
  +\sum_i\theta_i(x_i)
  +\sum_i\theta_{Yi}(y,x_i) \notag\\
  &+\sum_{(i,j)\in E_X}\theta_{ij}(x_i,x_j)
  +\sum_{(i,j)\in E_Y}\theta_{Yij}(y,x_i,x_j)
  -\log Z.
  \label{eq:full_model}
\end{align}
The chosen edge set \(E_X\) can be sparse, and the selected triples in \(E_Y\)
need not coincide with the past--past edges.

Once the past is known, the predictor is
\begin{equation}
  q(y\mid x)
  =
  \frac{\exp S_y(x)}
       {\sum_{y'\in\mcA}\exp S_{y'}(x)},
  \label{eq:conditional_score}
\end{equation}
with score
\begin{equation}
  S_y(x)
  =
  \theta_Y(y)
  +\sum_i\theta_{Yi}(y,x_i)
  +\sum_{(i,j)\in E_Y}\theta_{Yij}(y,x_i,x_j).
  \label{eq:score}
\end{equation}
The purely past-dependent terms in \eqref{eq:full_model} cancel from the
conditional probability.  This is algebraically true, but it does \emph{not}
mean that the measured past--past marginals are irrelevant.

\subsection{How past--past pairs influence a prediction}

The potential \(\theta_{ij}(x_i,x_j)\) does not appear directly in
\eqref{eq:score}.  However, it is coupled to the other potentials through all
the simultaneous moment-matching equations \eqref{eq:moment_match}.  Changing
the measured marginal \(\hatp_{ij}\) changes the fitted
\(\theta_{ij}\), and generally forces changes in
\(\theta_{Yi}\), \(\theta_{Yj}\), and any nearby triple potentials to keep the
target--past marginals and triples correct.  These changed target-dependent
potentials alter the prediction.

The role of \(\hatp_{ij}\) is to say how much of the information in \(X_i\) and
\(X_j\) is duplicated.  A model that knows they nearly always agree should not
treat two matching observations as two independent pieces of evidence.

\subsection{A concrete binary example}

Let \(Y,X_1,X_2\) be binary, each marginally uniform.  Suppose
\[
  \Pp(X_1=Y)=\Pp(X_2=Y)=0.9.
\]
If the only constraints are the two target--past pairs, maximum entropy gives
conditional independence \(X_1\ind X_2\mid Y\).  Therefore
\[
  \Pp(Y=1\mid X_1=1,X_2=1)
  =
  \frac{0.9^2}{0.9^2+0.1^2}
  \approx 0.9878.
\]
This is the familiar naive-Bayes multiplication of evidence.

Now add the measured past--past constraint
\[
  \Pp(X_1=X_2)=1.
\]
The only compatible support has \(X_1=X_2\).  Together with the two
target--past pair marginals, it implies
\[
  \Pp(Y=1,X_1=X_2=1)=0.45,\qquad
  \Pp(Y=0,X_1=X_2=1)=0.05.
\]
Thus
\[
  \Pp(Y=1\mid X_1=1,X_2=1)=\frac{0.45}{0.45+0.05}=0.9.
\]
The second observation contributes no new information because it is a perfect
duplicate of the first.  In a potential representation this hard dependence is
expressed by \(\theta_{12}(0,1)=\theta_{12}(1,0)=-\infty\), or by a large
negative finite value after smoothing.  The factor cancels from the final
conditional score, but it changes which joint distributions are feasible and
therefore changes the fitted target-dependent potentials.

\subsection{What measured triples add}

Suppose \(\hatp_{Yij}\) is available for a selected pair \((i,j)\).  The
potential \(\theta_{Yij}(y,x_i,x_j)\) captures a genuinely three-variable
effect.  It may express that the pair \((X_i,X_j)\) has a particular meaning for
the next token that neither token has separately.

If \(Y,X_i,X_j\) were the only variables, this is completely explicit:
\[
  q(y\mid x_i,x_j)
  =
  \frac{\hatp_{Yij}(y,x_i,x_j)}
       {\sum_{y'}\hatp_{Yij}(y',x_i,x_j)}.
\]
With other variables present, the triple potential is fitted together with all
other factors; one cannot simply multiply raw triple tables without correcting
for their overlaps.

\section{Exact easy cases}

\subsection{A tree of pair marginals}

Suppose the constraint graph consists only of one-variable and pairwise
marginals, and its graph is a tree \(T=(V,E)\).  If the measured edge marginals
are mutually consistent, the unique tree distribution is
\begin{equation}
  q_T(z)
  =
  \prod_{v\in V}\hatp_v(z_v)
  \prod_{(u,v)\in E}
  \frac{\hatp_{uv}(z_u,z_v)}
       {\hatp_u(z_u)\hatp_v(z_v)}.
  \label{eq:tree_factor}
\end{equation}
Equivalently,
\[
  q_T(z)
  =
  \frac{\prod_{(u,v)\in E}\hatp_{uv}(z_u,z_v)}
       {\prod_{v\in V}\hatp_v(z_v)^{\deg(v)-1}}.
\]
It matches every specified node and edge marginal.  It is also the
maximum-entropy extension of those tree constraints.

Formula \eqref{eq:tree_factor} gives the potentials directly:
\[
  \theta_v(z_v)=\log\hatp_v(z_v),\qquad
  \theta_{uv}(z_u,z_v)
  =
  \log\frac{\hatp_{uv}(z_u,z_v)}
                {\hatp_u(z_u)\hatp_v(z_v)}.
\]
This is exceptional: in a tree, the potentials are available directly from the
marginals, without iterative fitting.

If \(Y\) has neighbors \(N(Y)\), then with all past variables observed,
\[
  q_T(y\mid x)
  \propto
  \hatp_Y(y)^{1-\lvert N(Y)\rvert}
  \prod_{i\in N(Y)}\hatp_{Yi}(y,x_i).
\]
Only the Markov blanket \(N(Y)\) enters directly.  More distant past variables
affect \(Y\) only through those neighbors; once the neighbors are observed,
they carry no additional information in a tree model.

\subsection{Junction trees with pair and triple cliques}

The same simplification holds for a decomposable or junction-tree model.  Let
\(\mathcal K\) be clique marginals arranged in a clique tree, and let
\(\mathcal S\) be its separator marginals.  If they agree on intersections,
\begin{equation}
  q(z)
  =
  \frac{\prod_{C\in\mathcal K}\hatp_C(z_C)}
       {\prod_{S\in\mathcal S}\hatp_S(z_S)}.
  \label{eq:junction_tree}
\end{equation}
Repeated separators are included once for each clique-tree edge.

For example, if the cliques are
\[
  C_1=\{Y,X_1,X_2\},\qquad
  C_2=\{X_2,X_3\},
\]
the separator is \(\{X_2\}\) and
\[
  q(y,x_1,x_2,x_3)
  =
  \frac{\hatp_{Y12}(y,x_1,x_2)\hatp_{23}(x_2,x_3)}
       {\hatp_2(x_2)}.
\]
Conditioning on all the past gives
\[
  q(y\mid x_1,x_2,x_3)
  =
  q(y\mid x_1,x_2)
  =
  \frac{\hatp_{Y12}(y,x_1,x_2)}
       {\sum_{y'}\hatp_{Y12}(y',x_1,x_2)}.
\]
The \(X_2,X_3\) factor is important for the full joint distribution and for
missing-data inference, but does not directly change the prediction once
\(X_1,X_2,X_3\) are all observed.

\begin{keyidea}
Trees and junction trees are the computationally clean regime.  Their clique
marginals determine the full joint distribution by products and separator
divisions.  No global partition-function calculation is needed.
\end{keyidea}

\section{Fitting a general loopy model}

When the selected clique hypergraph has cycles, raw clique marginals cannot
usually be combined by multiplying them.  Their overlaps would be counted
incorrectly.  The potentials must be calibrated.

\subsection{Iterative proportional fitting}

If the full joint table were small enough to store, iterative proportional
fitting (IPF) would be the most direct algorithm.  Initialize with a positive
reference distribution \(q^{(0)}\).  For each clique \(C\), update
\begin{equation}
  q^{(r+1)}(z)
  =
  q^{(r)}(z)
  \frac{\hatp_C(z_C)}{q^{(r)}_C(z_C)}.
  \label{eq:ipf}
\end{equation}
The update makes the \(C\)-marginal correct, but may perturb previous
marginals; cycling through cliques converges under the usual positivity and
consistency conditions.  IPF is also an alternating KL-projection algorithm.

For a large alphabet, \eqref{eq:ipf} is a conceptual reference rather than an
implementation: a full joint table is impossible.  It nevertheless explains
where the potentials come from.  Each clique update adds
\[
  \log\hatp_C(z_C)-\log q_C^{(r)}(z_C)
\]
to the clique potential.  The needed quantity is the current model marginal
\(q_C^{(r)}\), which must be computed by inference in the factor graph.

\subsection{Gradient fitting}

One can instead maximize the dual objective \eqref{eq:dual}.  The gradient
\eqref{eq:gradient} has a simple interpretation:
\[
  \text{update}
  =
  \text{measured clique frequency}
  -
  \text{model clique frequency}.
\]
This is the usual fitting rule for an exponential family.  It is also the form
of the gradient used by log-linear models and conditional random fields.
Generalized iterative scaling is a classical alternative for such models
\cite{darroch_ratcliff}.

\section{Message passing and the meaning of tree-like}

Write each factor as
\[
  \phi_C(z_C)=\exp\{\theta_C(z_C)\},
  \qquad
  q(z)\propto\prod_C\phi_C(z_C).
\]
The sum-product algorithm passes messages between variable and factor nodes.
For a variable \(v\) and neighboring factor \(C\),
\begin{align}
  m_{v\to C}(z_v)
  &=
  \prod_{D\ni v,\ D\neq C}m_{D\to v}(z_v), \label{eq:var_message}\\
  m_{C\to v}(z_v)
  &=
  \sum_{z_{C\setminus v}}
  \phi_C(z_C)
  \prod_{u\in C\setminus v}m_{u\to C}(z_u).
  \label{eq:factor_message}
\end{align}

\subsection{When message passing is exact}

If the factor graph is a tree, one inward and one outward pass computes exact
marginals.  If the graph is not a tree but has small treewidth, it can be
converted to a junction tree and exact message passing remains possible.  The
cost is exponential in the largest clique size after this conversion, not in
the total number of variables.

Thus tree-like has two different meanings:
\begin{itemize}
  \item A \emph{tree} has treewidth one and admits exact, inexpensive
  sum-product messages.
  \item A sparse graph with cycles may still have low treewidth, in which case
  exact junction-tree inference is possible, but its cost grows as
  \(\lvert\mcA\rvert^{w+1}\), where \(w\) is the treewidth.
\end{itemize}

\subsection{Loopy belief propagation}

On a graph with larger cycles or high treewidth, the same message equations can
be iterated.  This is loopy belief propagation.  It is often useful, but it is
an approximation: it can fail to converge, can have multiple fixed points, and
its resulting clique marginals need not be exactly globally consistent.  It is
reasonable as a numerical approximation after validation on small cases, not as
an automatic replacement for the exact tree calculation.  The factor-graph and
sum-product formalism is described clearly in \cite{kschischang}.

\section{The large-alphabet obstacle}

Even a single dense pair table over \(10^5\) token values has \(10^{10}\)
entries, and a dense triple table has \(10^{15}\).  Therefore an implementation
cannot literally keep all potentials or all measured marginals.

Several representations are compatible with the graphical-model framework:
\begin{itemize}
  \item \textbf{Sparse exceptions plus a backoff model.}
  Store only frequent pairs or triples and use a smoothed lower-order
  distribution for all other configurations.
  \item \textbf{Token classes or a hierarchical vocabulary.}
  Model coarse token types first and refine within a class.
  \item \textbf{Low-rank nonnegative factors.}
  Represent a pair factor as
  \[
    \phi_{ij}(a,b)
    \approx
    \sum_{r=1}^R A_{ir}(a)B_{jr}(b).
  \]
  Then a factor-to-variable message can be computed in roughly
  \(O(R\lvert\mcA\rvert)\), rather than \(O(\lvert\mcA\rvert^2)\).
  Triple factors admit analogous nonnegative tensor decompositions.
  \item \textbf{Shared embeddings or hashed features.}
  Replace a separate parameter for every token identity by a lower-dimensional
  representation.  This loses the exact marginal-matching interpretation, but
  can be statistically far more sensible.
\end{itemize}

For a compressor, the fallback must give every symbol positive probability.
Otherwise a previously unseen configuration yields an infinite code length.

\section{Noisy marginals and regularization}
\label{sec:noisy}

Empirical marginals should not normally be imposed as exact constraints.
Rare triples are especially noisy, and different empirical tables may not be
perfectly consistent.  Three standard remedies are:
\begin{enumerate}
  \item \textbf{Smooth the measured tables.}  Mix each table with a lower-order
  backoff distribution before treating it as a target.
  \item \textbf{Use soft constraints.}  Instead of \(q_C=\hatp_C\), minimize
  \[
    \KL(q\Vert q_0)
    +
    \sum_{C\in\mcC}\lambda_C
    D_C(q_C,\hatp_C),
  \]
  where \(D_C\) penalizes disagreement and \(\lambda_C\) reflects confidence
  in the measurement.
  \item \textbf{Regularize the potentials.}  Fit the factor model by a
  likelihood or prequential code-length objective plus, for example,
  \(-\lambda\lVert\theta\rVert_2^2/2\).  This shrinks poorly supported
  interactions toward zero.
\end{enumerate}

The third form is often called conditional maximum entropy when the model is
fit directly as \(q_\theta(Y\mid \X)\).

\section{Conditional modelling versus full-joint modelling}

\subsection{The important simplification when all context tokens are known}

Suppose that at time \(t\) the compressor knows every selected past token
\[
  x=(x_1,\ldots,x_k).
\]
This is the usual causal prediction setting: encoder and decoder have both
decoded the same prefix.  There is then no need to infer the values of the
\(X_i\).  A full joint factor that depends only on \(x\) is a number and
cancels when the distribution is normalized over \(y\).  The online task is
therefore only
\[
  q(y\mid x)
  =
  \frac{\exp\{s(y,x)\}}
       {\sum_{y'\in\mcA}\exp\{s(y',x)\}},
  \label{eq:observed_context_softmax}
\]
for a score \(s(y,x)\).  In particular, ordinary sum-product messages over
the past variables are \emph{not} needed during this prediction step.

This distinction is useful:
\begin{itemize}
  \item A \emph{full-joint model} uses \(p(X_i,X_j)\) and other past-only
  statistics to construct one coherent distribution over both \(Y\) and
  \(X\).  It may require iterative fitting or message passing while the model
  is being built.
  \item A \emph{conditional model} builds only the score in
  \eqref{eq:observed_context_softmax}.  It never represents the probability of
  the observed context \(x\), so it does not need past-only potentials at
  prediction time.
\end{itemize}

\subsection{The exact easy case: a decomposable joint model}

First suppose the selected pair and triple tables form a consistent junction
tree.  The joint distribution is then the clique-over-separator product in
\eqref{eq:junction_tree}.  Let \(\mathcal K_Y\) denote its cliques containing
\(Y\), and \(\mathcal S_Y\) its separators containing \(Y\).  Since all other
factors are independent of \(y\), conditioning on the observed \(x\) gives
\begin{equation}
  q(y\mid x)
  \propto
  \frac{
    \displaystyle\prod_{C\in\mathcal K_Y}
    \hatp_C\bigl(y,x_{C\setminus\{Y\}}\bigr)}
  {
    \displaystyle\prod_{S\in\mathcal S_Y}
    \hatp_S\bigl(y,x_{S\setminus\{Y\}}\bigr)
  }.
  \label{eq:junction_conditional}
\end{equation}
The normalization in \eqref{eq:observed_context_softmax} completes the
calculation.  Thus a decomposable graph does not require iterative fitting or
message passing for a fully observed context: look up the relevant clique
values for each candidate \(y\), form the score, and normalize.

For example, if the only relevant clique is
\(\{Y,X_i,X_j\}\), then
\[
  q(y\mid x_i,x_j)
  =
  \frac{\hatp_{Yij}(y,x_i,x_j)}
       {\sum_{y'}\hatp_{Yij}(y',x_i,x_j)}
  =
  \frac{\hatp_{Yij}(y,x_i,x_j)}
       {\hatp_{ij}(x_i,x_j)}.
  \label{eq:triple_conditional}
\]
Here the past--past pair \(\hatp_{ij}\) enters directly: it is the
normalizing denominator that turns a joint triple table into a conditional
next-token distribution.

\subsection{A direct conditional predictor from pair and triple tables}

The decomposable case is exact but restrictive.  A more flexible construction
uses the measured tables as local pieces of predictive evidence.  Define,
after causal smoothing,
\begin{align}
  r_i(y\mid a)
  &=
  \frac{\hatp_{Yi}(y,a)}{\hatp_i(a)}, \label{eq:pair_predictor}\\
  r_{ij}(y\mid a,b)
  &=
  \frac{\hatp_{Yij}(y,a,b)}{\hatp_{ij}(a,b)}.
  \label{eq:triple_predictor}
\end{align}
The first is the predictor from one lag; the second is the predictor from a
pair of lags.  Their difference on the logarithmic scale is the conditional
interaction
\begin{equation}
  I_{ij}(y;a,b)
  =
  \log r_{ij}(y\mid a,b)
  -\log r_i(y\mid a)
  -\log r_j(y\mid b)
  +\log q_0(y).
  \label{eq:conditional_interaction}
\end{equation}
It measures the part of the pair prediction not already accounted for by the
two individual predictors, relative to a baseline \(q_0\).

One concrete conditional model is
\begin{equation}
  s(y,x)
  =
  \log q_0(y)
  +\sum_i w_i\log\frac{r_i(y\mid x_i)}{q_0(y)}
  +\sum_{(i,j)\in E_Y}w_{ij}I_{ij}(y;x_i,x_j),
  \label{eq:conditional_composite_score}
\end{equation}
followed by the softmax in \eqref{eq:observed_context_softmax}.  The weights
can be fixed, selected by a held-out prequential code length, or learned
online.  Formula \eqref{eq:conditional_composite_score} avoids the common
mistake of treating correlated lags as independent evidence.  For two lags,
with \(w_i=w_j=w_{ij}=1\), the score reduces exactly to
\(\log r_{ij}(y\mid x_i,x_j)\): the triple term corrects the double counting
in the two pair terms.

This construction clarifies two distinct roles of
\(\hatp_{ij}(x_i,x_j)\):
\begin{enumerate}
  \item It is needed explicitly in \eqref{eq:triple_predictor}.  A triple
  joint frequency must be divided by the observed pair frequency to become a
  conditional prediction.
  \item It indicates how often the two context values co-occur and therefore
  how reliable the corresponding triple estimate is.  It can set smoothing
  strength, decide whether an interaction is stored, or regularize
  \(w_{ij}\).
\end{enumerate}

If no triple \(\hatp_{Yij}\) is available, the past--past pair
\(\hatp_{ij}\) alone does \emph{not} determine how the two lags jointly
predict \(Y\).  It says that the lags are correlated, but it does not reveal
whether their correlation is predictive of the next token.  At that point one
needs an additional assumption: a downweighted product of pair predictors, a
full-joint maximum-entropy completion, or a learned conditional model.

\subsection{Fitting a conditional model from causal records}

If the original causal records \((y_t,x_t)\) are retained, one need not first
turn them into fixed tables.  Fit directly
\begin{equation}
  q_\theta(y\mid x)
  =
  \softmax_y\left[
    b(y)+\sum_i f_i(y,x_i)+\sum_{(i,j)\in E_Y}f_{ij}(y,x_i,x_j)
  \right].
  \label{eq:conditional_model}
\end{equation}
For example, the features can be the log ratios in
\eqref{eq:pair_predictor}--\eqref{eq:conditional_interaction}, sparse table
entries, or shared low-dimensional embeddings.  With an \(\ell_2\) penalty,
maximize
\[
  \mathcal L(\theta)
  =
  \sum_t\log q_\theta(y_t\mid x_t)
  -\frac{\lambda}{2}\lVert\theta\rVert_2^2.
  \label{eq:conditional_likelihood}
\]
For a pair feature \(\theta_i(c,a)\), its gradient is
\[
  \frac{\partial\mathcal L}{\partial\theta_i(c,a)}
  =
  \sum_t\1\{x_{i,t}=a\}
  \left[\1\{y_t=c\}-q_\theta(c\mid x_t)\right]
  -\lambda\theta_i(c,a),
  \label{eq:conditional_gradient}
\]
and an interaction feature has the same form with
\(\1\{x_{i,t}=a,x_{j,t}=b\}\) in front.

This optimization can be done by stochastic gradient descent, online mirror
descent, FTRL, or another standard conditional-model optimizer.  It does use
iterations, but these iterations do \emph{not} compute a global graphical-model
partition function.  Each update only evaluates the softmax over possible
next tokens for an observed context.  For a fixed finite feature map, the
regularized negative conditional log likelihood is convex.

Past--past pairs do not appear as a separate factor in
\eqref{eq:conditional_model}; all context values are already known.  They can
still be used to choose which interactions to include, to set their
regularization, and to make the interaction features in
\eqref{eq:conditional_interaction} well defined.  This is the appropriate
role for them in a purely conditional implementation.

\subsection{When full-joint fitting is still useful}

There are good reasons to fit the full-joint maximum-entropy model even when
the final context is observed:
\begin{itemize}
  \item The available information may be only overlapping aggregate
  pair/triple tables rather than the original records.  Full-joint
  moment-matching uses these tables directly and enforces their consistency.
  \item One may want calibrated probabilities for the context itself, to fill
  in missing lags, to sample sequences, or to use latent features.
  \item The past--past statistics may be part of the modelling claim, rather
  than merely a device for choosing conditional features.
\end{itemize}

In this route, IPF or gradient fitting may need model clique marginals.  On a
tree or junction tree they are available exactly.  On a graph with cycles,
one may use exact variable elimination when the treewidth is small, or an
approximation such as loopy belief propagation otherwise.  Crucially, this
work belongs to fitting or recalibrating the model.  Once the model has been
fitted and all \(x_i\) are observed, online prediction again reduces to
evaluating the \(y\)-dependent factors and normalizing over \(y\).

\subsection{When message passing is genuinely needed at prediction time}

Messages return only when something in the context has not been observed.
If a subset \(X_{\mathrm{miss}}\) is missing or latent, then
\[
  q(y\mid x_{\mathrm{obs}})
  =
  \sum_{x_{\mathrm{miss}}}
  q(y,x_{\mathrm{miss}}\mid x_{\mathrm{obs}}),
\]
and sum-product or junction-tree messages perform this summation.  The same is
true if the predictor contains hidden states, uncertain tokenization, or
mixtures whose gate depends on unobserved variables.  With a complete decoded
past, none of these sums is present.

\subsection{The remaining large-alphabet cost}

Knowing all of \(x\) removes inference over the past, but a dense softmax in
\eqref{eq:observed_context_softmax} still costs \(O(\lvert\mcA\rvert)\) per
token.  Message passing does not remove this cost: it comes from representing
the next-token distribution itself.  A particularly useful exact construction
is a full-support baseline with sparse conditional corrections.  Write
\[
  s(y,x)=\log q_0(y)+\Delta(y,x),
\]
and let \(C(x)\) be the candidate set on which \(\Delta(y,x)\ne 0\).  Then
\begin{equation}
  \sum_y q_0(y)e^{\Delta(y,x)}
  =
  1+\sum_{y\in C(x)}q_0(y)
     \left(e^{\Delta(y,x)}-1\right).
  \label{eq:sparse_normalizer}
\end{equation}
Thus the normalizer is exact in \(O(\lvert C(x)\rvert)\) time rather than
\(O(\lvert\mcA\rvert)\).  The candidate set can be the union of tokens with
stored pair or triple corrections for the observed lag values, plus a small
set of high-probability baseline tokens.

If the correction is dense, exact compression needs a different structured
output representation, such as a hierarchical vocabulary, a class-based
model, or a low-rank factorization that supports exact normalization and
cumulative probabilities.  Sampling-based softmax approximations are useful
for training, but the final encoder and decoder must use the same normalized,
full-support distribution.

\subsection{Practical consequence for the intended compressor}

For the setting described here, the most direct initial implementation is a
causal conditional predictor: maintain smoothed pair and selected triple
counts, construct the score in \eqref{eq:conditional_composite_score}, and
learn or validate its weights by prequential code length.  Use a full-support
baseline \(q_0\), and restrict expensive triple corrections to frequent,
well-supported context pairs.  A tree or junction-tree version is the best
exact reference implementation.  A loopy full-joint fit is useful as an
offline calibration experiment, but it is not required in the per-token
prediction loop when the entire context is known.

\section{Bayesian averaging and relative entropy}

The specified marginals do not themselves choose either a single completion or
a probability distribution over completions.  Maximum or relative entropy is a
\emph{point-completion} rule.  Bayesian averaging is a different rule: it
keeps uncertainty about the completion and averages the predictions.
Neither is intrinsically more correct; the useful choice depends on the
available data, the desired computation, and the modelling assumptions one is
willing to make.

\subsection{What exactly is being averaged?}

Let \(q_\theta\) be a family of joint distributions and let
\(\pi(d\theta\mid D)\) be a posterior distribution after data \(D\).  If the
past \(x\) is treated as an externally given covariate, the Bayesian predictor
is
\begin{equation}
  q_{\mathrm B}(y\mid x,D)
  =
  \int q_\theta(y\mid x)\,\pi(d\theta\mid D).
  \label{eq:bayes_conditional}
\end{equation}
If \(q_\theta\) is instead a model for the joint stream and the newly observed
past \(x\) is informative about \(\theta\), the posterior predictive
conditional is
\begin{equation}
  q_{\mathrm B}(y\mid x,D)
  =
  \frac{\displaystyle\int q_\theta(y,x)\,\pi(d\theta\mid D)}
       {\displaystyle\int\sum_{y'}q_\theta(y',x)\,\pi(d\theta\mid D)}.
  \label{eq:bayes_joint_predictive}
\end{equation}
Equation \eqref{eq:bayes_joint_predictive} first averages joint
distributions and then conditions.  It is not generally the same as simply
averaging \(q_\theta(y\mid x)\) with unchanged posterior weights: models that
make the observed context \(x\) more probable receive more weight.

One can state the idea of averaging directly on the set of completions.  Given
a prior \(\pi_0(dq)\) on the full joint simplex, a formal hard-constraint
posterior is
\[
  \pi(dq\mid \hatp)
  \propto
  \mathbf{1}\{q\in\mathcal F\}\,\pi_0(dq).
\]
There is no canonical uniform average until \(\pi_0\) has been chosen.
The flat Dirichlet prior is one possible convention, but it is still a
substantive choice and it lives in a simplex with
\(\lvert\mcA\rvert^{k+1}\) cells.

\subsection{The global Dirichlet prior: natural but far too large}

For a finite joint alphabet \(\mcZ\), the most familiar conjugate choice is
\[
  q\sim\operatorname{Dir}\{\alpha q_0(z):z\in\mcZ\},
  \qquad \alpha>0,
\]
where \(q_0\) is a full-support baseline distribution and \(\alpha\) is its
strength in pseudo-observations.  If complete joint records have counts
\(n(z)\), then
\[
  q\mid D
  \sim
  \operatorname{Dir}\{\alpha q_0(z)+n(z):z\in\mcZ\},
  \qquad
  \Pp(Z_{\mathrm{new}}=z\mid D)
  =
  \frac{\alpha q_0(z)+n(z)}{\alpha+n}.
  \label{eq:global_dirichlet}
\]
Thus the conjugate prior and posterior predictive are completely explicit.
They are unusable here because \(\lvert\mcZ\rvert\) grows as
\(\lvert\mcA\rvert^{k+1}\).  This example nevertheless identifies what a
useful structured prior should preserve: a baseline distribution, a
pseudo-count strength, and local updates from data.

\subsection{An exactly conjugate Bayesian tree model}

For a fixed tree, there is a practical structured replacement for
\eqref{eq:global_dirichlet}.  Choose a root \(r\), direct every edge away from
it, and write
\begin{equation}
  q_\theta(z)
  =
  \theta_r(z_r)
  \prod_{v\ne r}
  \theta_{v\mid \operatorname{pa}(v)}
     \bigl(z_v\mid z_{\operatorname{pa}(v)}\bigr).
  \label{eq:directed_tree}
\end{equation}
Put independent Dirichlet priors on the root table and on each row of each
conditional table:
\[
  \theta_r\sim\operatorname{Dir}(\alpha_r),
  \qquad
  \theta_{v\mid u=a}\sim\operatorname{Dir}(\alpha_{v\mid a})
  \quad\text{for every directed edge }u\to v.
\]
For complete records, let \(n_r(b)\) count \(Z_r=b\), and let
\(n_{v\mid a}(b)\) count the records with \(Z_u=a,Z_v=b\).  Conjugacy gives
\[
  \theta_r\mid D\sim\operatorname{Dir}\{\alpha_r(b)+n_r(b)\}_b,
  \qquad
  \theta_{v\mid u=a}\mid D
  \sim
  \operatorname{Dir}\{\alpha_{v\mid a}(b)+n_{v\mid a}(b)\}_b.
\]
Consequently the posterior-mean local probabilities are
\begin{align}
  \bar\theta_r(b)
  &=
  \frac{\alpha_r(b)+n_r(b)}
       {\sum_{b'}\alpha_r(b')+n_r(b')}, \label{eq:tree_root_predictive}\\
  \bar\theta_{v\mid a}(b)
  &=
  \frac{\alpha_{v\mid a}(b)+n_{v\mid a}(b)}
       {\sum_{b'}\alpha_{v\mid a}(b')+n_{v\mid a}(b')}.
  \label{eq:tree_edge_predictive}
\end{align}
The posterior factors remain independent, so the posterior predictive joint
is obtained simply by substituting these smoothed local tables into
\eqref{eq:directed_tree}.  The next-token probability is then
\(\bar q(y,x\mid D)/\sum_{y'}\bar q(y',x\mid D)\).  This is a true Bayesian
average, not merely a plug-in fit.  Choosing
\(\alpha_{v\mid a}=\tau_{v\mid a}b_{v\mid a}\) makes \(b_{v\mid a}\) a
backoff distribution and \(\tau_{v\mid a}\) its effective sample size.

The tree prior is computationally easy because it stores only local
conditional tables and needs no partition function.  In a very large alphabet,
the local tables still require sparse storage, backoff, token classes, or
shared parameters; conjugacy does not by itself cure the large-alphabet
problem.

\subsection{Decomposable graphs: the hyper-Dirichlet extension}

For a decomposable undirected graph or junction tree, the corresponding
conjugate family is the \emph{hyper-Dirichlet} family
\cite{dawid_lauritzen}.  It is a distribution over compatible clique and
separator probabilities, concentrated on joint distributions that are Markov
with respect to the chosen graph.  With complete records it stays in the same
family after observing data.

Concretely, choose a clique tree.  Its prior pseudo-count tables must agree
under marginalization on every separator.  Observing a record adds one count
to every clique configuration and to every separator configuration selected by
that record.  The posterior hyperparameters are therefore prior compatible
counts plus observed compatible counts.  A perfect elimination order converts
the decomposable graph to a directed acyclic factorization, where the same
calculation is visible as ordinary Dirichlet updates of conditional tables.
This gives exact posterior prediction and exact message passing whenever the
cliques are small enough.

\subsection{General loopy models: formal conjugacy is not cheap}

For the log-linear model in \eqref{eq:full_model}, complete records have
likelihood
\[
  p(D\mid\theta)
  \propto
  \exp\{\langle T(D),\theta\rangle-n\log Z(\theta)\},
\]
where \(T(D)\) contains the measured feature counts.  Every regular exponential
family has a formal Diaconis--Ylvisaker conjugate family
\cite{diaconis_ylvisaker},
\[
  \pi(\theta\mid\eta,\nu)
  \propto
  \exp\{\langle\eta,\theta\rangle-\nu\log Z(\theta)\}.
  \label{eq:dy_prior}
\]
After \(n\) records the parameters update to
\((\eta+T(D),\nu+n)\).  This is mathematically natural, but it is usually
not computationally convenient: evaluating \(Z(\theta)\), its gradient, and
the normalizing constant of \eqref{eq:dy_prior} requires the same difficult
inference that obstructs fitting the loopy model.  Independent Dirichlet
priors on its overlapping factor tables are \emph{not} conjugate in general.

Posterior prediction for a loopy model therefore normally uses an
approximation: Markov-chain Monte Carlo, variational inference, a Laplace
approximation around a regularized fit, or sometimes loopy message passing.
The tree and low-treewidth cases are special precisely because this global
inference disappears.

\subsection{If only marginal tables have been retained}

The clean conjugate updates above use the original complete records.  If only
the aggregate estimates
\(\hatp_{Yi},\hatp_{ij},\hatp_{Yij}\) have been retained, the tables are
overlapping summaries of the same data.  Their count noise is correlated.  It
is therefore incorrect to multiply independent multinomial likelihoods for
each pair and triple table: this counts the same observations several times.

There are three principled options:
\begin{enumerate}
  \item Retain or reconstruct the original causal records and use their joint
  likelihood.  This is the cleanest route for Bayesian updating.
  \item Use a composite likelihood made from selected marginal tables, but
  choose its weights by held-out code length or by an explicit effective-sample
  size calculation.  It is an approximation, not exact Bayes.
  \item Approximate the vector of empirical margins by a joint Gaussian
  distribution,
  \[
    \hat m\mid\theta
    \approx
    \mathcal N\bigl(m_\theta,\Sigma_\theta/n\bigr),
  \]
  where \(m_\theta\) contains the model pair and triple marginals and
  \(\Sigma_\theta\) contains their shared-sample covariance.  This respects
  the fact that overlapping tables are dependent, but computing
  \(m_\theta\) and \(\Sigma_\theta\) again requires inference in the graph.
\end{enumerate}

\subsection{Relative entropy as a point decision}

Maximum entropy uses a uniform reference distribution.  For prediction, it is
usually preferable to begin with a baseline \(q_0\) and make the smallest
change needed to respect the selected information:
\begin{equation}
  q^\star
  =
  \arg\min_{q\in\mathcal F}
  \KL(q\Vert q_0).
  \label{eq:relative_entropy}
\end{equation}
The solution is
\[
  q^\star(z)
  \propto
  q_0(z)
  \exp\left[\sum_{C\in\mcC}\theta_C(z_C)\right].
\]
It is a point decision, not generally the same as the Bayesian posterior mean
predictor.  One Bayesian interpretation is that \(q^\star\) is the mode of an
entropy-centred prior restricted to the constraint set, for example a density
proportional to \(\exp\{-\tau\KL(q\Vert q_0)\}\).  The strength \(\tau\)
controls how strongly the prior favours the baseline.

The baseline can be a unigram model, a conventional context model, a cache
model, or eventually a language model.  The graphical factors then correct the
baseline only where the measured data support a correction.

\section{Uncertainty about which interactions to trust}

The previous constructions assumed that a fixed collection of pair and triple
statistics had already been chosen.  In practice, one may have many candidate
lags and interactions, with only some worth using.  It is natural to place a
prior on this selection rather than commit to one graph before prediction.

\subsection{Choose interaction terms, not raw marginals}

The elementary objects should be predictive effects, not arbitrary raw
marginal tables.  A triple table \(\hatp_{Yij}\) already contains its
\(Y,X_i\), \(Y,X_j\), and \(X_i,X_j\) projections.  Selecting that triple
while declaring one of its projections absent is therefore not a meaningful
independent choice.

A useful hierarchy is:
\begin{itemize}
  \item a main effect from one lag, represented by
  \(\log\{r_i(y\mid x_i)/q_0(y)\}\);
  \item a genuine pair-of-lags effect, represented by the conditional
  interaction \(I_{ij}(y;x_i,x_j)\) in
  \eqref{eq:conditional_interaction}.
\end{itemize}
Introduce one inclusion indicator \(z_e\in\{0,1\}\) for each such effect and,
as a first prior,
\[
  \Pp(z)
  =
  \prod_{e=1}^m\rho_e^{z_e}(1-\rho_e)^{1-z_e}.
  \label{eq:independent_structure_prior}
\]
The probabilities \(\rho_e\) can depend on the order of the interaction, its
causal count, or a shared Beta--Bernoulli sparsity prior.  A hierarchical prior
can require an interaction to be included only when its corresponding
lower-order effects are available.

\subsection{Literal averaging over maximum-entropy completions}

For a selected set \(z\), let \(\mathcal F_z\) impose exactly the corresponding
chosen marginal constraints and define the relative-entropy completion
\[
  q_z
  =
  \arg\min_{q\in\mathcal F_z}\KL(q\Vert q_0).
\]
The literal Bayesian model average is
\begin{equation}
  \bar q(y\mid x,D)
  =
  \sum_{z\in\{0,1\}^m}
  \Pp(z\mid D)\,q_z(y\mid x).
  \label{eq:literal_bma}
\end{equation}
The prior in \eqref{eq:independent_structure_prior} is only the starting
point.  A Bayesian weight should be
\(\Pp(z\mid D)\propto \Pp(z)\Pp(D\mid z)\), or, for a compressor, can be
replaced by a weight based on past prequential code length.

Equation \eqref{eq:literal_bma} is generally harder than fitting one model.
There are \(2^m\) terms, and the potentials in \(q_z\) are normally
\emph{refitted} when a constraint is removed or added.  Consequently the
normalization constants and potentials do not factorize over \(e\).  Each
individual prediction is simple when all \(x_i\) are observed, but summing the
exponentially many predictions is not.

\subsection{Soft reliability: an exact elimination of the subset sum}

There is a different Bayesian formulation in which averaging is genuinely
simpler.  Let \(D_e(q_e,\hatp_e)\) measure disagreement between a model
marginal and an estimated marginal.  Interpret \(z_e=1\) as saying that the
estimate is trustworthy, and use the soft-constraint model
\begin{equation}
  \pi(q,z\mid\hatp)
  \propto
  \pi_0(q)
  \prod_{e=1}^m
  \left[
    \rho_e\exp\{-\lambda_eD_e(q_e,\hatp_e)\}
  \right]^{z_e}
  (1-\rho_e)^{1-z_e}.
  \label{eq:soft_reliability_joint}
\end{equation}
Summing every \(z_e\) independently is now exact:
\begin{equation}
  \pi(q\mid\hatp)
  \propto
  \pi_0(q)
  \prod_{e=1}^m
  \left[
    (1-\rho_e)
    +
    \rho_e\exp\{-\lambda_eD_e(q_e,\hatp_e)\}
  \right].
  \label{eq:soft_reliability_marginal}
\end{equation}
Thus the exponential combinatorial sum has disappeared.  The price is that
this is not the same object as \eqref{eq:literal_bma}: it is a posterior over
joint distributions under uncertain \emph{constraint reliability}, rather than
an average of separately refitted hard-constraint models.

The induced penalty for one marginal is
\[
  \Psi_e(d)
  =
  -\log\left[
    (1-\rho_e)+\rho_e e^{-\lambda_e d}
  \right].
\]
It is near zero for a well-matched marginal and saturates at
\(-\log(1-\rho_e)\) for a badly mismatched one.  This saturation is attractive:
a dubious table can effectively switch itself off rather than forcing the
entire model to distort itself.  Inference over \(q\) may still be difficult
on a loopy graph, but no explicit enumeration of subsets remains.

\subsection{A fast conditional surrogate: annealed optional evidence}

When the context \(x\) is fully observed, let \(\delta_e(y,x)\) be a fixed
conditional effect, such as a main effect or an interaction in
\eqref{eq:conditional_composite_score}, and define
\[
  q_z(y\mid x)
  =
  \frac{
    q_0(y)\exp\{\sum_e z_e\delta_e(y,x)\}
  }{
    \sum_{y'}q_0(y')\exp\{\sum_e z_e\delta_e(y',x)\}
  }.
\]
The exact average \(\E_z q_z(y\mid x)\) is again an average of ratios and does
not factorize.  However, averaging the \emph{unnormalized} evidence does:
\begin{equation}
  q_{\mathrm{ann}}(y\mid x)
  \propto
  q_0(y)
  \prod_{e=1}^m
  \left[
    (1-\rho_e)+\rho_e e^{\delta_e(y,x)}
  \right].
  \label{eq:annealed_predictor}
\end{equation}
This annealed predictor is not identical to literal Bayesian model averaging,
because an expectation of a normalized distribution is not generally a
normalized expectation.  It is nevertheless a coherent, cheaply normalized
conditional model with soft effective potentials
\[
  \delta_e
  \longmapsto
  \log\left[(1-\rho_e)+\rho_e e^{\delta_e}\right].
\]
For a small effect,
\[
  \log\left[(1-\rho_e)+\rho_e e^{\delta_e}\right]
  =
  \rho_e\delta_e
  +
  \frac{\rho_e(1-\rho_e)}{2}\delta_e^2
  +O(\delta_e^3).
\]
To first order, uncertain evidence is simply tempered by \(\rho_e\).  If the
effects are stored sparsely, the exact sparse normalizer
\eqref{eq:sparse_normalizer} applies directly to
\eqref{eq:annealed_predictor}.

\subsection{The tree exception: exact averaging over many structures}

There is one important case where a very large model average is exactly
tractable.  Restrict the candidate pairwise structures to spanning trees and
use a product-form prior
\[
  \Pp(T)
  \propto
  \prod_{e\in T}w_e.
\]
The edges are not independent, because they must form a tree, but the weight
factorizes over selected edges.  If \(L\) is the weighted graph Laplacian,
\[
  L_{ij}=-w_{ij}\ (i\ne j),\qquad
  L_{ii}=\sum_{j\ne i}w_{ij},
\]
then the matrix-tree theorem gives
\begin{equation}
  \sum_{T\ \mathrm{spanning\ tree}}\prod_{e\in T}w_e
  =
  \det L^{(r)},
  \label{eq:matrix_tree}
\end{equation}
where \(L^{(r)}\) is any one row and its corresponding column removed.

With compatible local Dirichlet priors, complete records update the local
tables and the posterior edge weights.  For a proposed new configuration, its
local predictive factors modify the \(w_e\), so the tree average can again be
computed by a determinant.  This gives exact Bayesian prediction while
averaging over all spanning-tree structures in polynomial time
\cite{meila_jaakkola}.  The method is a clean reference for pairwise
dependence, but arbitrary triple interactions or unrestricted loopy graphs
lose this determinant simplification.

\subsection{Recommended use}

For an initial compressor, the most promising practical choice is the
annealed conditional predictor \eqref{eq:annealed_predictor}: define sparse
main and interaction effects from causal tables, use \(\rho_e\) as a
reliability or inclusion strength, and tune these strengths by prequential
code length.  The soft-reliability posterior
\eqref{eq:soft_reliability_marginal} is the more principled route when one
wants uncertainty over noisy marginal constraints themselves.  A Bayesian
mixture of trees provides an exact benchmark.  Literal averaging of arbitrary
maximum-entropy subsets is best reserved for small validation problems or
Monte Carlo experiments.

\section{A concrete route for a causal compressor}

The following route retains the graphical-model interpretation while staying
computationally realistic.

\begin{enumerate}
  \item Select a modest set of lag variables \(X_i\).  Include a safe base
  predictor with full support.
  \item Estimate smoothed target--past pair statistics for candidate lags.
  Use them to identify individual predictive signals.
  \item Estimate selected past--past pair statistics.  Use these to identify
  redundancy and to propose a sparse tree or low-treewidth graph.
  \item Add target--past--past triples only where held-out or prequential code
  length shows a genuine interaction beyond the pair model.
  \item Prefer a tree or junction-tree structure when possible.  Then compute
  the predictor directly from clique and separator marginals.
  \item If a loopy model is necessary, use sparse or low-rank factors; validate
  loopy message passing against exact small instances.
  \item Update all estimates causally from previously decoded tokens only, with
  identical finite-precision rules in encoder and decoder.
\end{enumerate}

\begin{keyidea}
The criterion for adding a pair, triple, or edge should be reduction in
prequential code length after including the cost of the additional parameters.
This is the compression form of model selection.
\end{keyidea}

\section{Conclusion}

The correct general object is not a single pairwise model, nor a single
triple model.  It is a constraint hypergraph containing whichever pair and
triple marginals are trustworthy.  Maximum or relative entropy selects one
coherent joint model, while Bayesian averaging keeps a distribution over
plausible models and averages their predictions.  The latter is exactly
tractable for a fixed tree with local Dirichlet priors, and for decomposable
graphs through the hyper-Dirichlet construction.  When it is unclear which
interactions deserve trust, soft reliability variables remove the exponential
subset sum exactly; mixtures over all tree structures provide a separate
determinant-based exact benchmark.

In trees and junction trees, the construction is explicit and exact.  In
general loopy graphs, potentials must be fitted to match all constraints, and
Bayesian inference as well as message passing becomes approximate unless the
graph has small treewidth.  If only overlapping marginal tables have been
saved, their shared sampling noise must be modelled or approximated; they
cannot be treated as independent observations.
For a very large token alphabet, the essential practical issue is not only
graph structure but the representation of each factor: sparse corrections,
backoff distributions, classes, and low-rank factors are necessary.

\begin{thebibliography}{9}

\bibitem{wainwright_jordan}
M. J. Wainwright and M. I. Jordan.
\newblock Graphical Models, Exponential Families, and Variational Inference.
\newblock \emph{Foundations and Trends in Machine Learning}, 1(1--2):1--305,
2008.
\newblock \url{https://www.cs.columbia.edu/~blei/fogm/2020F/readings/WainwrightJordan2008.pdf}.

\bibitem{chow_liu}
C. K. Chow and C. N. Liu.
\newblock Approximating Discrete Probability Distributions with Dependence Trees.
\newblock \emph{IEEE Transactions on Information Theory}, 14(3):462--467, 1968.
\newblock \url{https://cs.nyu.edu/~roweis/csc2515-2006/readings/chowliu.pdf}.

\bibitem{darroch_ratcliff}
J. N. Darroch and D. Ratcliff.
\newblock Generalized Iterative Scaling for Log-Linear Models.
\newblock \emph{The Annals of Mathematical Statistics}, 43(5):1470--1480,
1972.
\newblock \url{https://doi.org/10.1214/aoms/1177692379}.

\bibitem{kschischang}
F. R. Kschischang, B. J. Frey, and H.-A. Loeliger.
\newblock Factor Graphs and the Sum-Product Algorithm.
\newblock \emph{IEEE Transactions on Information Theory}, 47(2):498--519,
2001.
\newblock \url{https://www.isiweb.ee.ethz.ch/papers/arch/aloe-2001-1.pdf}.

\bibitem{diaconis_ylvisaker}
P. Diaconis and D. Ylvisaker.
\newblock Conjugate Priors for Exponential Families.
\newblock \emph{The Annals of Statistics}, 7(2):269--281, 1979.
\newblock \url{https://doi.org/10.1214/aos/1176344611}.

\bibitem{dawid_lauritzen}
A. P. Dawid and S. L. Lauritzen.
\newblock Hyper Markov Laws in the Statistical Analysis of Decomposable
Graphical Models.
\newblock \emph{The Annals of Statistics}, 21(3):1272--1317, 1993.
\newblock \url{https://doi.org/10.1214/aos/1176349260}.

\bibitem{meila_jaakkola}
M. Meil\u{a} and T. S. Jaakkola.
\newblock Tractable Bayesian Learning of Tree Belief Networks.
\newblock Technical Report CMU-RI-TR-00-15, Carnegie Mellon University, 2000.
\newblock \url{https://sites.stat.washington.edu/mmp/Papers/prior-final-form.pdf}.

\end{thebibliography}

\end{document}
